% Deck 2 of 4 of AI for Biodiversity: A Primer.
% Build: latexmk -pdf main.tex   (pdflatex + biber)
\documentclass[11pt,aspectratio=169]{beamer}
% ---------------------------------------------------------------------------
% preamble.tex --- shared Beamer style for the four workshop decks
%
% Loaded by each deck with % ---------------------------------------------------------------------------
% preamble.tex --- shared Beamer style for the four workshop decks
%
% Loaded by each deck with % ---------------------------------------------------------------------------
% preamble.tex --- shared Beamer style for the four workshop decks
%
% Loaded by each deck with \input{../common/preamble}. Same fonts, colours
% and bibliography settings as the book (book/preamble.tex): Latin Modern,
% black on white, one grey for code, one accent colour for what is untrusted
% or wrong. Built with pdflatex + biber, like the book.
% ---------------------------------------------------------------------------

\usepackage[T1]{fontenc}
\usepackage[utf8]{inputenc}
\usepackage{lmodern}
\usepackage{microtype}
\usepackage{booktabs}
\usepackage{tabularx}
\usepackage{listings}
\usepackage{xcolor}
\usepackage{csquotes}
\usepackage{graphicx}

% Figures: the book's figures as the 300 dpi PNGs the book itself includes
% (book/figures/png, tracked in git and current; figures/pdf is not rebuilt
% after every change), and the first pages of the papers
% (common/screenshots). Both are found by name without an extension.
\graphicspath{{../../book/figures/png/}{../common/screenshots/}}
\DeclareGraphicsExtensions{.png,.pdf}

% --- look ------------------------------------------------------------------
\usetheme{default}
\usefonttheme{serif}
\definecolor{accent}{HTML}{9B2D20}
\definecolor{codebg}{gray}{0.96}
\definecolor{rule}{gray}{0.75}
\definecolor{dim}{gray}{0.45}
\setbeamercolor{structure}{fg=black}
\setbeamercolor{normal text}{fg=black,bg=white}
\setbeamercolor{alerted text}{fg=accent}
\setbeamercolor{section in toc}{fg=black}
\setbeamercolor{footline}{fg=dim}
\setbeamerfont{title}{series=\bfseries,size=\Large}
\setbeamerfont{subtitle}{series=\mdseries,size=\normalsize}
\setbeamerfont{frametitle}{series=\bfseries,size=\large}
\setbeamerfont{footline}{size=\scriptsize}
\setbeamertemplate{navigation symbols}{}
\setbeamertemplate{itemize items}[circle]
\setbeamertemplate{enumerate items}[default]
\setbeamertemplate{caption}{\scriptsize\insertcaption}
\setbeamertemplate{frametitle}{%
  \vskip0.6em\usebeamerfont{frametitle}\insertframetitle\par
  \vskip0.2em{\color{rule}\hrule height 0.4pt}\vskip0.4em}
\setbeamertemplate{footline}{%
  \hbox to \paperwidth{\hspace{1.2em}\usebeamerfont{footline}\usebeamercolor[fg]{footline}%
    \insertshorttitle\ --- \insertsection\hfill
    \insertframenumber/\inserttotalframenumber\hspace{1.2em}}\vskip0.6em}
\setbeamersize{text margin left=1.4em,text margin right=1.4em}
% An outline frame at the start of each section, the current one in black.
\AtBeginSection[]{\begin{frame}{Outline}\tableofcontents[currentsection]\end{frame}}
% A frame that breaks over pages is numbered only from its second page.
\setbeamertemplate{frametitle continuation}[from second]

% --- code listings (as in the book) ----------------------------------------
\lstset{
  basicstyle=\ttfamily\footnotesize,
  backgroundcolor=\color{codebg},
  frame=none,
  breaklines=true,
  breakatwhitespace=true,
  columns=fullflexible,
  keepspaces=true,
  upquote=true,
  showstringspaces=false,
  xleftmargin=0.6em,
  framexleftmargin=0.6em,
  aboveskip=0.5em,
  belowskip=0.5em,
  literate={~}{{\textasciitilde}}1
           {>}{{\textgreater}}1
           {<}{{\textless}}1
           {ö}{{\"o}}1 {ü}{{\"u}}1 {é}{{\'e}}1,
}
\lstdefinestyle{shell}{language=bash}

% --- bibliography (same settings as the book) ------------------------------
\usepackage[
  backend=biber,
  style=authoryear,
  sorting=nyt,
  maxbibnames=6,
  minbibnames=3,
  maxcitenames=2,
  giveninits=true,
  uniquename=init,
  doi=true,
  url=true,
  isbn=false,
  eprint=true,
]{biblatex}
\addbibresource{../../book/references.bib}
\renewcommand*{\bibfont}{\scriptsize}
\setlength{\bibitemsep}{0.25em}
\setbeamertemplate{bibliography item}{}   % no icon before each entry

% --- shorthand -------------------------------------------------------------
\newcommand{\code}[1]{\texttt{#1}}
\newcommand{\term}[1]{\emph{#1}}
% A key point, as in the book: no box.
\newcommand{\kp}[1]{\par\medskip{\bfseries Key point.}\enspace\emph{#1}\par}
% \fig[height]{name}: one of the book's figures, scaled to fit the frame.
\newcommand{\fig}[2][0.72\textheight]{%
  \centering\includegraphics[height=#1,width=\textwidth,keepaspectratio]{#2}\par}
% \screenshot{key}: the first page of a paper, with a thin frame.
\newcommand{\screenshot}[2][0.66\textheight]{%
  \setlength{\fboxsep}{0pt}\setlength{\fboxrule}{0.3pt}%
  \fbox{\includegraphics[height=#1,width=\linewidth,keepaspectratio]{#2}}}
% \paperslide{key}{text}: screenshot on the left, what the book takes from the
% paper on the right. Use inside a frame; add \source after it.
\newcommand{\paperslide}[2]{%
  \begin{columns}[T]
    \begin{column}{0.40\textwidth}\centering\screenshot{#1}\end{column}
    \begin{column}{0.58\textwidth}\small #2\end{column}
  \end{columns}}
% Notes in references.bib record how a web source was saved and dated; they
% belong in the reference list, not on a \source line, so cite commands
% (including \fullcite) leave them out.
\AtEveryCitekey{\clearfield{note}}
% \source{...}: a small grey line at the foot of a frame, for the source.
\newcommand{\source}[1]{\vskip0pt plus 1filll\par{\scriptsize\color{dim}#1\par}}
. Same fonts, colours
% and bibliography settings as the book (book/preamble.tex): Latin Modern,
% black on white, one grey for code, one accent colour for what is untrusted
% or wrong. Built with pdflatex + biber, like the book.
% ---------------------------------------------------------------------------

\usepackage[T1]{fontenc}
\usepackage[utf8]{inputenc}
\usepackage{lmodern}
\usepackage{microtype}
\usepackage{booktabs}
\usepackage{tabularx}
\usepackage{listings}
\usepackage{xcolor}
\usepackage{csquotes}
\usepackage{graphicx}

% Figures: the book's figures as the 300 dpi PNGs the book itself includes
% (book/figures/png, tracked in git and current; figures/pdf is not rebuilt
% after every change), and the first pages of the papers
% (common/screenshots). Both are found by name without an extension.
\graphicspath{{../../book/figures/png/}{../common/screenshots/}}
\DeclareGraphicsExtensions{.png,.pdf}

% --- look ------------------------------------------------------------------
\usetheme{default}
\usefonttheme{serif}
\definecolor{accent}{HTML}{9B2D20}
\definecolor{codebg}{gray}{0.96}
\definecolor{rule}{gray}{0.75}
\definecolor{dim}{gray}{0.45}
\setbeamercolor{structure}{fg=black}
\setbeamercolor{normal text}{fg=black,bg=white}
\setbeamercolor{alerted text}{fg=accent}
\setbeamercolor{section in toc}{fg=black}
\setbeamercolor{footline}{fg=dim}
\setbeamerfont{title}{series=\bfseries,size=\Large}
\setbeamerfont{subtitle}{series=\mdseries,size=\normalsize}
\setbeamerfont{frametitle}{series=\bfseries,size=\large}
\setbeamerfont{footline}{size=\scriptsize}
\setbeamertemplate{navigation symbols}{}
\setbeamertemplate{itemize items}[circle]
\setbeamertemplate{enumerate items}[default]
\setbeamertemplate{caption}{\scriptsize\insertcaption}
\setbeamertemplate{frametitle}{%
  \vskip0.6em\usebeamerfont{frametitle}\insertframetitle\par
  \vskip0.2em{\color{rule}\hrule height 0.4pt}\vskip0.4em}
\setbeamertemplate{footline}{%
  \hbox to \paperwidth{\hspace{1.2em}\usebeamerfont{footline}\usebeamercolor[fg]{footline}%
    \insertshorttitle\ --- \insertsection\hfill
    \insertframenumber/\inserttotalframenumber\hspace{1.2em}}\vskip0.6em}
\setbeamersize{text margin left=1.4em,text margin right=1.4em}
% An outline frame at the start of each section, the current one in black.
\AtBeginSection[]{\begin{frame}{Outline}\tableofcontents[currentsection]\end{frame}}
% A frame that breaks over pages is numbered only from its second page.
\setbeamertemplate{frametitle continuation}[from second]

% --- code listings (as in the book) ----------------------------------------
\lstset{
  basicstyle=\ttfamily\footnotesize,
  backgroundcolor=\color{codebg},
  frame=none,
  breaklines=true,
  breakatwhitespace=true,
  columns=fullflexible,
  keepspaces=true,
  upquote=true,
  showstringspaces=false,
  xleftmargin=0.6em,
  framexleftmargin=0.6em,
  aboveskip=0.5em,
  belowskip=0.5em,
  literate={~}{{\textasciitilde}}1
           {>}{{\textgreater}}1
           {<}{{\textless}}1
           {ö}{{\"o}}1 {ü}{{\"u}}1 {é}{{\'e}}1,
}
\lstdefinestyle{shell}{language=bash}

% --- bibliography (same settings as the book) ------------------------------
\usepackage[
  backend=biber,
  style=authoryear,
  sorting=nyt,
  maxbibnames=6,
  minbibnames=3,
  maxcitenames=2,
  giveninits=true,
  uniquename=init,
  doi=true,
  url=true,
  isbn=false,
  eprint=true,
]{biblatex}
\addbibresource{../../book/references.bib}
\renewcommand*{\bibfont}{\scriptsize}
\setlength{\bibitemsep}{0.25em}
\setbeamertemplate{bibliography item}{}   % no icon before each entry

% --- shorthand -------------------------------------------------------------
\newcommand{\code}[1]{\texttt{#1}}
\newcommand{\term}[1]{\emph{#1}}
% A key point, as in the book: no box.
\newcommand{\kp}[1]{\par\medskip{\bfseries Key point.}\enspace\emph{#1}\par}
% \fig[height]{name}: one of the book's figures, scaled to fit the frame.
\newcommand{\fig}[2][0.72\textheight]{%
  \centering\includegraphics[height=#1,width=\textwidth,keepaspectratio]{#2}\par}
% \screenshot{key}: the first page of a paper, with a thin frame.
\newcommand{\screenshot}[2][0.66\textheight]{%
  \setlength{\fboxsep}{0pt}\setlength{\fboxrule}{0.3pt}%
  \fbox{\includegraphics[height=#1,width=\linewidth,keepaspectratio]{#2}}}
% \paperslide{key}{text}: screenshot on the left, what the book takes from the
% paper on the right. Use inside a frame; add \source after it.
\newcommand{\paperslide}[2]{%
  \begin{columns}[T]
    \begin{column}{0.40\textwidth}\centering\screenshot{#1}\end{column}
    \begin{column}{0.58\textwidth}\small #2\end{column}
  \end{columns}}
% Notes in references.bib record how a web source was saved and dated; they
% belong in the reference list, not on a \source line, so cite commands
% (including \fullcite) leave them out.
\AtEveryCitekey{\clearfield{note}}
% \source{...}: a small grey line at the foot of a frame, for the source.
\newcommand{\source}[1]{\vskip0pt plus 1filll\par{\scriptsize\color{dim}#1\par}}
. Same fonts, colours
% and bibliography settings as the book (book/preamble.tex): Latin Modern,
% black on white, one grey for code, one accent colour for what is untrusted
% or wrong. Built with pdflatex + biber, like the book.
% ---------------------------------------------------------------------------

\usepackage[T1]{fontenc}
\usepackage[utf8]{inputenc}
\usepackage{lmodern}
\usepackage{microtype}
\usepackage{booktabs}
\usepackage{tabularx}
\usepackage{listings}
\usepackage{xcolor}
\usepackage{csquotes}
\usepackage{graphicx}

% Figures: the book's figures as the 300 dpi PNGs the book itself includes
% (book/figures/png, tracked in git and current; figures/pdf is not rebuilt
% after every change), and the first pages of the papers
% (common/screenshots). Both are found by name without an extension.
\graphicspath{{../../book/figures/png/}{../common/screenshots/}}
\DeclareGraphicsExtensions{.png,.pdf}

% --- look ------------------------------------------------------------------
\usetheme{default}
\usefonttheme{serif}
\definecolor{accent}{HTML}{9B2D20}
\definecolor{codebg}{gray}{0.96}
\definecolor{rule}{gray}{0.75}
\definecolor{dim}{gray}{0.45}
\setbeamercolor{structure}{fg=black}
\setbeamercolor{normal text}{fg=black,bg=white}
\setbeamercolor{alerted text}{fg=accent}
\setbeamercolor{section in toc}{fg=black}
\setbeamercolor{footline}{fg=dim}
\setbeamerfont{title}{series=\bfseries,size=\Large}
\setbeamerfont{subtitle}{series=\mdseries,size=\normalsize}
\setbeamerfont{frametitle}{series=\bfseries,size=\large}
\setbeamerfont{footline}{size=\scriptsize}
\setbeamertemplate{navigation symbols}{}
\setbeamertemplate{itemize items}[circle]
\setbeamertemplate{enumerate items}[default]
\setbeamertemplate{caption}{\scriptsize\insertcaption}
\setbeamertemplate{frametitle}{%
  \vskip0.6em\usebeamerfont{frametitle}\insertframetitle\par
  \vskip0.2em{\color{rule}\hrule height 0.4pt}\vskip0.4em}
\setbeamertemplate{footline}{%
  \hbox to \paperwidth{\hspace{1.2em}\usebeamerfont{footline}\usebeamercolor[fg]{footline}%
    \insertshorttitle\ --- \insertsection\hfill
    \insertframenumber/\inserttotalframenumber\hspace{1.2em}}\vskip0.6em}
\setbeamersize{text margin left=1.4em,text margin right=1.4em}
% An outline frame at the start of each section, the current one in black.
\AtBeginSection[]{\begin{frame}{Outline}\tableofcontents[currentsection]\end{frame}}
% A frame that breaks over pages is numbered only from its second page.
\setbeamertemplate{frametitle continuation}[from second]

% --- code listings (as in the book) ----------------------------------------
\lstset{
  basicstyle=\ttfamily\footnotesize,
  backgroundcolor=\color{codebg},
  frame=none,
  breaklines=true,
  breakatwhitespace=true,
  columns=fullflexible,
  keepspaces=true,
  upquote=true,
  showstringspaces=false,
  xleftmargin=0.6em,
  framexleftmargin=0.6em,
  aboveskip=0.5em,
  belowskip=0.5em,
  literate={~}{{\textasciitilde}}1
           {>}{{\textgreater}}1
           {<}{{\textless}}1
           {ö}{{\"o}}1 {ü}{{\"u}}1 {é}{{\'e}}1,
}
\lstdefinestyle{shell}{language=bash}

% --- bibliography (same settings as the book) ------------------------------
\usepackage[
  backend=biber,
  style=authoryear,
  sorting=nyt,
  maxbibnames=6,
  minbibnames=3,
  maxcitenames=2,
  giveninits=true,
  uniquename=init,
  doi=true,
  url=true,
  isbn=false,
  eprint=true,
]{biblatex}
\addbibresource{../../book/references.bib}
\renewcommand*{\bibfont}{\scriptsize}
\setlength{\bibitemsep}{0.25em}
\setbeamertemplate{bibliography item}{}   % no icon before each entry

% --- shorthand -------------------------------------------------------------
\newcommand{\code}[1]{\texttt{#1}}
\newcommand{\term}[1]{\emph{#1}}
% A key point, as in the book: no box.
\newcommand{\kp}[1]{\par\medskip{\bfseries Key point.}\enspace\emph{#1}\par}
% \fig[height]{name}: one of the book's figures, scaled to fit the frame.
\newcommand{\fig}[2][0.72\textheight]{%
  \centering\includegraphics[height=#1,width=\textwidth,keepaspectratio]{#2}\par}
% \screenshot{key}: the first page of a paper, with a thin frame.
\newcommand{\screenshot}[2][0.66\textheight]{%
  \setlength{\fboxsep}{0pt}\setlength{\fboxrule}{0.3pt}%
  \fbox{\includegraphics[height=#1,width=\linewidth,keepaspectratio]{#2}}}
% \paperslide{key}{text}: screenshot on the left, what the book takes from the
% paper on the right. Use inside a frame; add \source after it.
\newcommand{\paperslide}[2]{%
  \begin{columns}[T]
    \begin{column}{0.40\textwidth}\centering\screenshot{#1}\end{column}
    \begin{column}{0.58\textwidth}\small #2\end{column}
  \end{columns}}
% Notes in references.bib record how a web source was saved and dated; they
% belong in the reference list, not on a \source line, so cite commands
% (including \fullcite) leave them out.
\AtEveryCitekey{\clearfield{note}}
% \source{...}: a small grey line at the foot of a frame, for the source.
\newcommand{\source}[1]{\vskip0pt plus 1filll\par{\scriptsize\color{dim}#1\par}}

% Six of this deck's figures (ch11-two-layers, ch11-citation-check,
% ch12-where-data-goes, ch12-architecture, ch12-kybele, ch14-scoring-grid)
% have no PDF under book/figures/pdf/, which is untracked and predates the
% renumbering of 14 October 2026. The book's 300 dpi PNGs are the fallback:
% graphicx tries .pdf before .png in every directory, so a PDF is used
% wherever one exists. Remove this line once scripts/build_figures.sh has run.
% Note: beamer reads a brace group right after a frame title as the frame's
% subtitle, so \fig is wrapped in \begingroup ... \endgroup, never in braces.

\title[AI for Biodiversity 2]{AI for Biodiversity: A Primer}
\subtitle{2. Working with agents, and local-first (Chapters 7--14)}
\author{Savvas Paragkamian}
\institute{Hellenic Centre for Marine Research \\ KYBELE and ELIXIR-GR}
\date{HCMR workshop, November 2026}
\hypersetup{pdftitle={AI for Biodiversity: A Primer, deck 2},pdfauthor={Savvas Paragkamian}}

\begin{document}

\begin{frame}[plain]
  \titlepage
\end{frame}

\begin{frame}{Goals of this session}
  \begin{enumerate}
    \item Work with a coding agent the way a system administrator works:
      describe the result, set limits on what the agent can change, and
      check its work (Chapter 7).
    \item Keep an analysis reproducible when one of its tools may give a
      different answer each time it runs (Chapters 8--11).
    \item Decide what runs locally, on what hardware, and how to test a
      model on your own data before relying on it (Chapters 12--14).
  \end{enumerate}
  \medskip
  Running examples: the invented springtail survey of Chapter 10
  (\code{workshop/session}), and KYBELE as Chapter 12 describes it, a
  design rather than a running system. Chapters 10, 11 and 14 get one or
  two frames here; deck 3 walks through them.
\end{frame}

% ---------------------------------------------------------------------------
\section{The agentic coding workflow (Chapter 7)}

\begin{frame}{Vibe coding and CHOP}
  \begingroup\fig[0.34\textheight]{ch07-chop}\endgroup
  \begin{itemize}
    \item \term{Vibe coding} means accepting generated code without reading
      it. Fine for a throwaway plot; not for anything entering a manuscript,
      a database or someone else's pipeline.
    \item \term{Chat-oriented programming} (CHOP) is the useful alternative:
      you describe what you want, the system proposes code, you run and
      check it, and you refine.
    \item The difference is whether checking is in the loop. You are
      accountable for what you publish, whatever wrote it.
  \end{itemize}
  \source{Figure from Chapter 7 of the book.}
\end{frame}

\begin{frame}[fragile]{The project speaks through files}
  \code{README.md} explains the project to people and agents: if a new
  postdoc could not start from it, neither can an agent. \code{AGENTS.md}
  is the contract, and each line should close a failure you have seen.
\begin{lstlisting}
## Data
- data/raw/ is READ-ONLY. Derived data goes to data/processed/
  and must be regenerable by running scripts/ in order.
- Record the CRS of every spatial file; occurrence
  coordinates are WGS84. Never reproject silently.
- Match taxon names to a named backbone before any join,
  and record its version or download date.
  Never join on raw name strings.
# ...
\end{lstlisting}
  Cut: the Code section (R via \code{renv}, no global installs) and the
  Don't section (don't commit, don't touch \code{manuscript/}, don't add
  dependencies without asking).
  \source{Chapter 7, \emph{The project speaks through files}.}
\end{frame}

\begin{frame}{Review: six checks}
  \begin{enumerate}\small
    \item \textbf{Commit before.} A clean checkpoint makes everything
      reversible.
    \item \textbf{Read the whole diff.} If it is too big to read, the task
      was too big: revert and split it.
    \item \textbf{Test the edges:} empty input, one row, missing values,
      duplicate keys.
    \item \textbf{Count rows before and after every join.} Silently lost
      rows are, in the author's experience, the most common serious error
      in ecological data processing, and agents produce them readily.
    \item \textbf{Run it twice.} Different output means hidden state or an
      unset random seed.
    \item \textbf{Check the science.} The agent has no field experience and
      will not flag a biologically impossible result.
  \end{enumerate}
  \kp{Agents fail loudly on syntax and quietly on meaning. Spend your
    review time on meaning.}
  \source{Chapter 7, \emph{Review}.}
\end{frame}

\begin{frame}{Does it save time? Measure your own}
  \paperslide{becker2025metr}{%
    \begin{itemize}
      \item A randomised trial, 16 experienced open-source developers:
        tasks took 19 per cent longer with AI tools, yet the developers
        believed they had saved 20 per cent \parencite{becker2025metr}.
      \item That measured the tools of early 2025. METR's follow-up
        pointed to a speed-up, but METR judged the estimates unreliable:
        developers declined to take part, or held back tasks, rather than
        work without AI \parencite{becker2026uplift}.
      \item A meta-analysis of 27 effect sizes: a moderate overall gain,
        smaller in open-source and enterprise work
        \parencite{maier2026metaanalysis}.
      \item Measure your own time before assuming either.
    \end{itemize}}
  \source{\fullcite{becker2025metr}}
\end{frame}

\begin{frame}{Analysis with agents: keep the judgement calls}
  Statistical analysis is where code that runs is furthest from code that
  is right. In a trial at Anthropic, 50 practitioners analysed data with an
  agent secretly set to favour one of two reasonable options at each
  decision; they went with the agent's side 86 per cent of the time, where
  chance would give half \parencite{herdagdelen2026aligned}. Five habits:
  \begin{itemize}\small
    \item \textbf{Fix the plan before the session:} question, response,
      predictors, model family and exclusion rules, as a preregistration
      would, and give the agent that file.
    \item \textbf{Choose the model yourself.} If the agent proposes another,
      ask it to make the case for each option before you decide.
    \item \textbf{Log every model fitted,} so the report can say how many
      were tried. Make this a rule in \code{AGENTS.md}.
    \item \textbf{Check the diagnostics yourself:} residuals, and simulated
      data with known parameters.
    \item \textbf{Report how the result changes} under other reasonable
      choices.
  \end{itemize}
  \source{Chapter 7, \emph{Analysis with agents}.}
\end{frame}

% ---------------------------------------------------------------------------
\section{The software stack and reproducibility (Chapters 8--9)}

\begin{frame}{The layers, and four ways to talk to a system}
  \begingroup\fig[0.6\textheight]{ch08-layers}\endgroup
  For reproducibility the command line wins: a command is a record of what
  you did, a series of clicks is not. The book's harnesses run in the
  terminal because a session there leaves a record that can be replayed.
  \source{Figure from Chapter 8 of the book.}
\end{frame}

\begin{frame}{Local model tooling: two settings cause most failures}
  \begingroup\fig[0.48\textheight]{ch08-local-tooling}\endgroup
  \begin{itemize}
    \item Tell the harness that the model supports tools, and enlarge the
      context: the runtime's default is sized for chat, and an agent loop
      re-sends the whole session each turn.
    \item Ollama also serves models hosted on its own servers through the
      same command line. Check that a model is local before trusting it
      with data.
  \end{itemize}
  \source{Figure from Chapter 8 of the book. Current defaults and versions: Appendix D.}
\end{frame}

\begin{frame}{Where a language model breaks reproducibility}
  \begingroup\fig[0.45\textheight]{ch09-variation}\endgroup
  \begin{itemize}
    \item Sampling is random: temperature zero reduces the variation but
      does not remove it \parencite{coqueret2026randomness}.
      \textcite{keck2026interactions} reran 100 paragraphs 100 times each
      at temperature zero, and the runs did not always agree.
    \item A hosted model with a fixed name is not a fixed thing. A local
      model file with a known checksum does not change: the strongest
      argument for open weights in research.
  \end{itemize}
  \source{Figure from Chapter 9 of the book.}
\end{frame}

\begin{frame}{Three places for a model in an analysis}
  \begingroup\fig[0.64\textheight]{ch09-chain-of-inference}\endgroup
  \kp{Generated code is reproducible; generated answers are not. Move every
    model-dependent step as far upstream as it will go, and keep the
    deterministic artefact.}
  \source{Figure from Chapter 9 of the book.}
\end{frame}

\begin{frame}[fragile]{Log model calls like instrument readings}
  \code{logs/2026-10-15\_taxon-matching.md}, as the book prints it:
\begin{lstlisting}
Model: qwen3-coder:30b, Q4_K_M, sha256:4f2a...c91
Runtime: ollama 0.31.1   Harness: opencode 1.17.11   Context: 64K
Temperature: 0   Seed: 42   Fresh session   Rules: AGENTS.md @ 7c1e2d0
Prompt: Write an R script matching data/raw/species_list.csv against
  the GBIF backbone; keep unmatched rows as status="NO_MATCH".
Reference: GBIF Backbone Taxonomy as of 2026-10-15, DOI in README.md
Output: scripts/03_match_taxa.R (commit a3f9c21)
Review: 1,847 rows in, 1,847 out. 23 NO_MATCH checked by hand:
  19 typos fixed in a lookup table, 4 names not in the backbone
  referred to a specialist.
\end{lstlisting}
  Each line records something that can change the result: the runtime and
  harness versions, the quantisation, the context length, the rules file,
  and the backbone the names were matched against. Keep it in plain
  Markdown beside the analysis, and commit it with the code.
  \source{Chapter 9, \emph{Log model calls like instrument readings}.}
\end{frame}

% ---------------------------------------------------------------------------
\section{An agent session and a weekly review (Chapters 10--11)}

\begin{frame}{Rules as advice, rules as configuration}
  \begin{columns}[T]
    \begin{column}{0.52\textwidth}
      \fig[0.6\textheight]{ch10-three-layers}
    \end{column}
    \begin{column}{0.46\textwidth}\small
      \begin{itemize}
        \item \code{AGENTS.md} is advice: the model can misread or
          ignore it.
        \item The harness configuration binds only the harness's own
          tools; a script the agent writes and runs is an ordinary program
          with your permissions.
        \item A sandbox, where there is one, confines every process the
          agent starts; by default the project folder stays writable.
        \item File permissions: the operating system refuses the write,
          whoever asks.
      \end{itemize}
    \end{column}
  \end{columns}
  \source{Figure from Chapter 10 of the book. Each layer has a weakness the
    others cover; deck 3 runs the session itself step by step.}
\end{frame}

\begin{frame}{A weekly literature review: two layers, never mixed}
  \begingroup\fig[0.45\textheight]{ch11-two-layers}\endgroup
  \begin{itemize}
    \item Scripts in R with fixed queries, run by a clock: fetch five
      sources, merge duplicates, check every identifier, write the digest,
      the BibTeX file and the log. No model is involved.
    \item A Claude Code skill built on KYBELE's rules reads the week's
      files and writes notes in a file that says a model wrote it. A
      reference enters a bibliography only through a script that checked it.
  \end{itemize}
  \source{Figure from Chapter 11 of the book.}
\end{frame}

\begin{frame}{Plain files: a person, a script and an agent can all read them}
  \paperslide{piskala2026files}{%
    \begin{itemize}
      \item The project speaks through files: \code{README.md},
        \code{AGENTS.md}, skills, and memory you can read and edit.
        Memory you cannot see is a reproducibility problem (Chapter 7).
      \item Everything the weekly review produces is a plain file, which a
        person, a script and an agent can all read without an intermediary
        (Chapter 11; \textcite{piskala2026files}).
      \item The week's log is its instrument reading: the run time, the
        commit of the configuration, the checksum of every raw file, every
        query, every count and every failure, so that a week can be rerun
        and compared.
    \end{itemize}}
  \source{\fullcite{piskala2026files}}
\end{frame}

% ---------------------------------------------------------------------------
\section{The local-first approach (Chapter 12)}

\begin{frame}{Local-first software: you own your data}
  \paperslide{kleppmann2019localfirst}{%
    \begin{itemize}
      \item What software must do if users are to own their data: be
        fast, work across devices and offline, support collaboration,
        last, be private by default, and leave the user in control
        \parencite{kleppmann2019localfirst}.
      \item It predates the AI boom, but its questions also help assess
        how we use models and who controls the data.
      \item For a public institution: could you carry on if a provider
        changed its terms, raised prices or left your market? Can you audit
        what you depend on? Will your infrastructure outlive the contract?
    \end{itemize}}
  \source{\fullcite{kleppmann2019localfirst}}
\end{frame}

\begin{frame}{The case for local, and what it costs}
  \begin{columns}[T]
    \begin{column}{0.49\textwidth}\small
      \begin{itemize}
        \item \textbf{Nothing leaves the machine:} unpublished records,
          sensitive localities, embargoed sequences, drafts. Check it, don't
          assume it: watch the network traffic once.
        \item \textbf{No cost per query:} the whole corpus, not a sample;
          jobs run overnight.
        \item \textbf{It works offline:} vessels, field stations, closed
          networks.
        \item \textbf{It lasts:} a model file on your disk does not change;
          research infrastructure has a ten-year horizon.
      \end{itemize}
    \end{column}
    \begin{column}{0.49\textwidth}\small
      \begin{itemize}
        \item \textbf{Capability:} hosted frontier models lead on hard
          reasoning, long documents and many-file changes; open-weight
          models matched the best closed score 3.5 months later on average
          \parencite{emberson2025openweightlag}.
        \item \textbf{Speed:} slow on a laptop.
        \item \textbf{Hardware and upkeep:} money up front, and time out of
          someone's research.
        \item \textbf{A dissent:} for someone starting out, local models
          are by and large not worth it; begin with a hosted model
          \parencite{wickham2026nobs}. A taxonomic test supports the
          caution \parencite{schmidtlebuhn2025morphology}.
      \end{itemize}
    \end{column}
  \end{columns}
  \source{Chapter 12, \emph{The case for local} and \emph{The costs}.}
\end{frame}

\begin{frame}{Where the data goes, and a group policy}
  \begingroup\fig[0.42\textheight]{ch12-where-data-goes}\endgroup
  \begin{itemize}\small
    \item \textbf{Always local:} unpublished data, sensitive localities,
      personal data, anything under embargo or a restrictive agreement.
    \item \textbf{Local by default:} bulk processing, extraction and
      classification over large collections.
    \item \textbf{Hosted is fine:} public data, published literature,
      general explanation, hard reasoning where the gap matters.
    \item \textbf{Hosted with care:} code generation; check what context is
      sent and whether your licences and policy allow it.
  \end{itemize}
  \source{Figure from Chapter 12 of the book.}
\end{frame}

\begin{frame}{A reference architecture for a group}
  \begingroup\fig[0.62\textheight]{ch12-architecture}\endgroup
  The interface is for talking, the models are for thinking, and plain text
  is for remembering. A small proxy that writes every exchange to a dated
  Markdown file gives the provenance record Chapter 9 asks for.
  \source{Figure from Chapter 12 of the book.}
\end{frame}

\begin{frame}{KYBELE: a design}
  \begingroup\fig[0.66\textheight]{ch12-kybele}\endgroup
  \source{Figure from Chapter 12 of the book. KYBELE is the author's project: a
    literature assistant for Collembola on the SIB Literature Services. Its
    serving layer runs, its orchestrator is designed but not built, and
    none of its tests has been run.}
\end{frame}

\begin{frame}{Testing KYBELE, as planned}
  A suite to rerun after every change of prompt, model or retrieval. None
  of it has been run yet:
  \begin{itemize}\small
    \item a gold set of expert Collembola questions with expected sources:
      did it find them?
    \item a grounding check: is each claim supported by what was retrieved?
    \item a citation audit: does the cited paper say this?
    \item a pushback test: does it hold a correct answer when challenged?
    \item the token cost of each answer, and a comparison with the previous
      version, so that replacing the model is accepted only if the gold-set
      results hold.
  \end{itemize}
  Evidence for the design comes from two tested systems of the same kind,
  over Reactome and over Conservation Evidence, where combined keyword and
  vector retrieval did best
  \parencite{mohammadi2025reacttome,iyer2025conservationevidence}.
  \kp{Put domain expertise into data, code and tests --- vocabulary, scope
    filter, gold set --- before putting it into model weights. Fine-tune
    last, and only keep a tuned model that passes the tests.}
  \source{Chapter 12, \emph{KYBELE: a design}.}
\end{frame}

% ---------------------------------------------------------------------------
\section{Hardware, and testing on your own task (Chapters 13--14)}

\begin{frame}{Memory, not compute}
  \begingroup\fig[0.37\textheight]{ch13-memory-bandwidth}\endgroup
  The context needs memory too: the KV cache grows with the context
  length, and at the lengths agents use it can rival the weights. A
  32\,GB machine comfortably runs a Q4 12B model with a 32K context; a 70B
  model wants 48\,GB or more.
  \kp{Capacity first: do the model and its context fit? Bandwidth sets how
    fast an answer streams; compute sets how long an agent waits before
    each answer. Buying for processing power alone is the common mistake;
    buying for memory alone is the mistake an agent user makes.}
  \source{Figure from Chapter 13 of the book.}
\end{frame}

\begin{frame}{A size guide}
  \begingroup\fig[0.4\textheight]{ch13-size-guide}\endgroup
  \begin{itemize}\small
    \item Classification needs 1--7B models; summarising and extraction
      7--14B; scoped coding 12--32B with at least a 32K context; agent work
      across several files 30--70B. Below that, expect frustration.
    \item A mixture of experts reads only its active parameters per token:
      30B-A3B needs the memory of a 30B model and decodes like a 3B one.
    \item \code{Q4\_K\_M} is the default quantisation worth starting from.
      Test the quantised model on your own documents before relying on it.
  \end{itemize}
  \source{Figure from Chapter 13 of the book. Unified memory and particular
    machines: Chapter 13 and Appendix D.}
\end{frame}

\begin{frame}{Test several things at once}
  \paperslide{carli2026karenina}{%
    \begin{itemize}
      \item Karenina, from EMBL-EBI, lets experts say in plain language
        what a sound answer needs and what behaviour to test, and turns
        that into runnable checks. A preprint, not yet peer reviewed
        \parencite{carli2026karenina}.
      \item On 144 drug-discovery questions, live access to Open Targets
        raised the average accuracy of seven models, yet 11 per cent of
        those answers were correct but unsupported by the retrieved
        evidence.
      \item When a user confidently pushed back with false information,
        models often abandoned correct answers: sycophancy, measured. A
        second model as a guard flagged almost all reversals.
    \end{itemize}}
  \source{\fullcite{carli2026karenina}}
\end{frame}

\begin{frame}{A procedure for ecological tasks}
  \begingroup\fig[0.58\textheight]{ch14-procedure}\endgroup
  Whether a model is good enough for your data is something to test on your
  data. The exercise in \code{workshop/testing} runs this procedure on 25
  invented field notes from the survey; deck 3 shows the runs.
  \source{Figure from Chapter 14 of the book: the author's proposal, drawing
    on the extraction and classification studies of Chapter 4.}
\end{frame}

% ---------------------------------------------------------------------------
\begin{frame}{Takeaways}
  \begin{itemize}
    \item Agents fail loudly on syntax and quietly on meaning. Spend your
      review time on meaning.
    \item Generated code is reproducible; generated answers are not. Keep
      the deterministic artefact, and log a model call like an instrument
      reading.
    \item The rules file is advice. The harness configuration, its sandbox
      and the file permissions are enforcement, each with gaps. Use them
      together.
    \item A reference enters a bibliography through a script that resolved
      it, never through a model.
    \item Put domain expertise into data, code and tests before putting it
      into model weights.
    \item Capacity first: do the model and its context fit? Then test the
      model on your own task, with a key, a baseline and a score.
  \end{itemize}
\end{frame}

\begin{frame}[allowframebreaks]{References}
  \printbibliography[heading=none]
\end{frame}

\end{document}
